% ============================================
% DOC-04
% Gerado de molambudos.md — build_miolo.py
% ============================================
\section*{\textbf{DOC-04} --- Relatório de Sessão I}
\markboth{DOC-04 --- Relatório de Sessão I}{}

\noindent \textit{Dr. Heitor Oliveira --- 14 de março de 1979}

\noindent \textbf{Dr. Heitor Oliveira --- 14 de março de 1979}

\vspace{0.5em}\noindent\rule{\textwidth}{0.5pt}\vspace{0.5em}

\noindent \textbf{Paciente:} Joaquim Antônio Correia --- Paciente 1.260
\textbf{Tipo:} Sessão de avaliação psiquiátrica inicial
\textbf{Duração:} 45 minutos

\vspace{0.5em}\noindent\rule{\textwidth}{0.5pt}\vspace{0.5em}

\noindent \textbf{Observações do Paciente:}

\noindent Joaquim é um homem magro, de estatura baixa, pele enrugada e olhar fixo. O olho direito difere significativamente do esquerdo --- apresenta coloração amarelo-ouvo incomum. Quando questionado, afirmou que \textquotedblleft{}um bicho\textquotedblright{} mudou a cor. Recusou-se a elaborar.

\noindent O paciente é articulado e inteligente --- muito mais do que o prontuário sugere. Demonstra capacidade de meta-cognição e análise psicológica (ex: classificou médicos pelo cheiro, distinguiu entre médicos novos e experientes). Não apresenta os sinais clássicos de deterioração cognitiva esperados para alguém com 62 anos de institucionalização.

\noindent Seu discurso é lúcido, mas monofocado no tema da \textquotedblleft{}fome\textquotedblright{} e da \textquotedblleft{}criatura\textquotedblright{}. Quando questionado sobre os homicídios, respondeu com naturalidade: \textquotedblleft{}a fome matou por mim\textquotedblright{\textit{. Não demonstra remorso, mas também não demonstra orgulho --- demonstra }aceitação}, como se os eventos fossem inevitáveis.

\noindent \textit{Nota pessoal (manuscrito, à margem):} \textquotedblleft{}A chave deste caso não está no delírio. Está no que veio antes --- no sertão, na seca, no pai morto, na vala. O delírio é só a fechadura. A história é a chave. --- H.O.\textquotedblright{}

\noindent \textbf{Impressão Diagnóstica (Preliminar):}

\noindent O paciente apresenta quadro compatível com Esquizofrenia Paranóide (CID-8 295.3), com delírio sistematizado envolvendo uma entidade que o habita. A estrutura do delírio é complexa e internamente coerente --- incomum para pacientes institucionalizados por décadas, que geralmente apresentam discurso fragmentado.

\noindent No entanto, há elementos que sugerem capacidade de insight preservado. Quando perguntei se ele achava que a criatura era real ou fruto da mente, ele respondeu: \textit{\textquotedblleft{}O senhor acha que a diferença importa?\textquotedblright{}}

\noindent \textbf{Plano Terapêutico:}

\noindent - Sessões semanais de avaliação
- Solicitar acesso ao diário do paciente (informação: o paciente mantém um diário há décadas; nunca permitiu que ninguém lesse)
- Manter medicação atual (Clorpromazina 100mg/dia)

\noindent \textbf{Nota Pessoal (não incluída no prontuário oficial):}

\noindent Há algo neste paciente que me incomoda. Não consigo identificar precisamente. Mas ao final da sessão, quando me levantei para sair, ele disse: \textit{\textquotedblleft{}O senhor vai sonhar comigo esta noite.\textquotedblright{}}

\noindent Eu sonhei.

\vspace{0.5em}\noindent\rule{\textwidth}{0.5pt}\vspace{0.5em}

\noindent \textit{Dr. Heitor Oliveira}
\textit{Psiquiatra --- Ala de Segurança Máxima}
\textit{Hospital Colônia de Barbacena}

\vspace{0.5em}
\noindent\textit{\textbf{↪ Rotas:} \rota{MEM-13} \rotasep \rota{DOC-05} \rotasep \rota{MEM-14} \rotasep \rota{LUC-01}}